% Options for packages loaded elsewhere
% Options for packages loaded elsewhere
/PassOptionsToPackage{unicode}{hyperref}
/PassOptionsToPackage{hyphens}{url}
/PassOptionsToPackage{dvipsnames,svgnames,x11names}{xcolor}
%
/documentclass[
  letterpaper,
  12pt]{article}
/usepackage{xcolor}
/usepackage[margin=1in]{geometry}
/usepackage{amsmath,amssymb}
/setcounter{secnumdepth}{-/maxdimen} % remove section numbering
/usepackage{iftex}
/ifPDFTeX
  /usepackage[T1]{fontenc}
  /usepackage[utf8]{inputenc}
  /usepackage{textcomp} % provide euro and other symbols
/else % if luatex or xetex
  /usepackage{unicode-math} % this also loads fontspec
  /defaultfontfeatures{Scale=MatchLowercase}
  /defaultfontfeatures[/rmfamily]{Ligatures=TeX,Scale=1}
/fi
/usepackage{lmodern}
/ifPDFTeX/else
  % xetex/luatex font selection
  /setmainfont[]{TeX Gyre Termes}
/fi
% Use upquote if available, for straight quotes in verbatim environments
/IfFileExists{upquote.sty}{/usepackage{upquote}}{}
/IfFileExists{microtype.sty}{% use microtype if available
  /usepackage[]{microtype}
  /UseMicrotypeSet[protrusion]{basicmath} % disable protrusion for tt fonts
}{}
% Make /paragraph and /subparagraph free-standing
/makeatletter
/ifx/paragraph/undefined/else
  /let/oldparagraph/paragraph
  /renewcommand{/paragraph}{
    /@ifstar
      /xxxParagraphStar
      /xxxParagraphNoStar
  }
  /newcommand{/xxxParagraphStar}[1]{/oldparagraph*{#1}/mbox{}}
  /newcommand{/xxxParagraphNoStar}[1]{/oldparagraph{#1}/mbox{}}
/fi
/ifx/subparagraph/undefined/else
  /let/oldsubparagraph/subparagraph
  /renewcommand{/subparagraph}{
    /@ifstar
      /xxxSubParagraphStar
      /xxxSubParagraphNoStar
  }
  /newcommand{/xxxSubParagraphStar}[1]{/oldsubparagraph*{#1}/mbox{}}
  /newcommand{/xxxSubParagraphNoStar}[1]{/oldsubparagraph{#1}/mbox{}}
/fi
/makeatother


/usepackage{longtable,booktabs,array}
/newcounter{none} % for unnumbered tables
/usepackage{calc} % for calculating minipage widths
% Correct order of tables after /paragraph or /subparagraph
/usepackage{etoolbox}
/makeatletter
/patchcmd/longtable{/par}{/if@noskipsec/mbox{}/fi/par}{}{}
/makeatother
% Allow footnotes in longtable head/foot
/IfFileExists{footnotehyper.sty}{/usepackage{footnotehyper}}{/usepackage{footnote}}
/makesavenoteenv{longtable}
/usepackage{graphicx}
/makeatletter
/newsavebox/pandoc@box
/newcommand*/pandocbounded[1]{% scales image to fit in text height/width
  /sbox/pandoc@box{#1}%
  /Gscale@div/@tempa{/textheight}{/dimexpr/ht/pandoc@box+/dp/pandoc@box/relax}%
  /Gscale@div/@tempb{/linewidth}{/wd/pandoc@box}%
  /ifdim/@tempb/p@</@tempa/p@/let/@tempa/@tempb/fi% select the smaller of both
  /ifdim/@tempa/p@</p@/scalebox{/@tempa}{/usebox/pandoc@box}%
  /else/usebox{/pandoc@box}%
  /fi%
}
% Set default figure placement to htbp
/def/fps@figure{htbp}
/makeatother





/setlength{/emergencystretch}{3em} % prevent overfull lines

/providecommand{/tightlist}{%
  /setlength{/itemsep}{0pt}/setlength{/parskip}{0pt}}






% APA Seventh Edition styling, in plain LaTeX.
%
% Everything here is ordinary article-class styling. There is no apa7: the
% title page, the abstract, the author note and the notes below floats all
% arrive as ordinary blocks from apaquarto's filters, and this file only says
% how they are set. Each command below is used by formatlatex.lua.

/usepackage{setspace}
/usepackage{fancyhdr}
/usepackage{titlesec}
/usepackage{ragged2e}
/usepackage{etoolbox}
% [H] placement, which floatsintext asks for
/usepackage{float}

% --- Page and paragraph -----------------------------------------------------
% APA: double spaced throughout, first line of every paragraph indented half an
% inch, no extra space between paragraphs, text flush left with a ragged right
% edge rather than justified.
% /RaggedRight sets /parindent to /RaggedRightParindent, so the indent has to
% be told to it before it is switched on. Setting /parindent afterwards instead
% leaves every paragraph flush left.
/setlength{/parindent}{0.5in}
/setlength{/parskip}{0pt}
/setlength{/RaggedRightParindent}{0.5in}
/RaggedRight

% setspace's own double spacing, which is what apa7 uses and which comes to
% about twice the font size. Asking for linestretch: 2 instead multiplies the
% normal baseline skip, which is already larger than the font size, and set the
% lines about a fifth further apart than apaquarto-pdf sets them.
/doublespacing

% --- Running head -----------------------------------------------------------
% APA seventh edition: the short title in capitals at the left of every page
% and the page number at the right. No rule under it.
/pagestyle{fancy}
/fancyhf{}
/renewcommand{/headrulewidth}{0pt}
/fancyhead[L]{/apashorttitle}
/fancyhead[R]{/thepage}
/newcommand{/apashorttitle}{}
/newcommand{/setapashorttitle}[1]{/renewcommand{/apashorttitle}{/MakeUppercase{#1}}}
/newcommand{/apaauthorline}{}
/newcommand{/setapaauthorline}[1]{/renewcommand{/apaauthorline}{/MakeUppercase{#1}}}

% The head a published article carries, which is the one apaquarto's typst
% format sets: the short title centred on a recto page and the authors'
% surnames centred on a verso one, both in capitals, with the page number in
% the outer corner. The head is smaller than the body and the page number is
% not, which is how the journals set it.
%
% Recto and verso only differ once the document is twoside, which
% formatlatex.lua asks for along with this.
/newcommand{/apaheadsize}{/fontsize{8}{10}/selectfont}
/newcommand{/apajournalhead}{%
  /fancyhf{}%
  /fancyhead[LE]{/thepage}%
  /fancyhead[RO]{/thepage}%
  /fancyhead[CE]{{/apaheadsize/apaauthorline}}%
  /fancyhead[CO]{{/apaheadsize/apashorttitle}}}

% Float placement in a published article.
%
% A float spanning both columns can only go at the top of a page, at the foot
% of one where stfloats allows it, or on a page of its own. Latex's own
% fractions are cautious enough that the two spanning floats in example.qmd
% went to a page of their own with no text on it at all, where apa7 and the
% typst format both set them at the head of a page of text. A spanning float
% may now take nine tenths of the top of a page, and a float page has to be
% nine tenths full before latex will make one.
/newcommand{/apajoufloats}{%
  /renewcommand{/dbltopfraction}{0.9}%
  /renewcommand{/dblfloatpagefraction}{0.9}%
  /renewcommand{/topfraction}{0.9}%
  /renewcommand{/bottomfraction}{0.5}%
  /renewcommand{/textfraction}{0.07}%
  /renewcommand{/floatpagefraction}{0.85}%
  /setcounter{dbltopnumber}{2}%
  /setcounter{topnumber}{3}%
  /setcounter{totalnumber}{5}}

% The gutter between the columns, which typst leaves at a twenty-fifth of the
% text width and latex at ten points.
/newcommand{/apajoucolumnsep}{/setlength{/columnsep}{0.04/textwidth}}

% The first page of a journal or document paper carries the masthead instead of
% a running head, so it is left plain.
/fancypagestyle{apafirstpage}{%
  /fancyhf{}%
  /renewcommand{/headrulewidth}{0pt}}

% Two columns, which is what a published article is set in. /twocolumn starts a
% page, so it is asked for at the start of the document rather than in the
% preamble, where it would be too early to take effect.
/newcommand{/apatwocolumn}{/AtBeginDocument{/twocolumn}}

% --- Journal front matter ---------------------------------------------------
% In a published article the title, the byline, the abstract and the lines
% under them span the page beneath the masthead rather than standing in the
% first column, and the author note goes to the foot of that column. The sizes
% are the ones apa7 gives its journal mode, which is where apaquarto's typst
% format takes them from as well: at a 10pt base the title is /LARGE and
% unemphasised, the byline /large, the affiliations the body size, and the
% abstract and everything with it /small.
/newcommand{/apajoutitle}[1]{%
  {/centering/fontsize{17.28}{20.74}/selectfont #1/par}}

% The byline, with the affiliations under it at the body size. formatlatex.lua
% switches to the smaller size after the first paragraph, which is the byline.
% The gaps between the parts of the front matter are measured against the
% typst masthead, which is what this page is meant to look like.
/newenvironment{apajoubyline}
  {/par/vspace{9pt}/centering/fontsize{12}{14}/selectfont}
  {/par}
/newcommand{/apajouaffiliationsize}{/fontsize{10}{12}/selectfont}

% The abstract, the impact statement and the lines under them: smaller, and in
% a block narrower than the masthead, centred under the byline. apa7 sets its
% journal abstract in a 4.6875in parbox, which is where the width comes from.
% The text inside is set flush left, as it is in the typst masthead.
/newlength{/apajounarrowwidth}
/setlength{/apajounarrowwidth}{4.6875in}
% parskip as well as leading: typst leaves a line and a bit between the
% keywords and the supplemental materials under them, where latex on its own
% would leave only the line.
/newenvironment{apajounarrow}
  {/par/vspace{7pt}/centering
   /begin{minipage}{/apajounarrowwidth}%
   /fontsize{9}{11}/selectfont
   /setlength{/parskip}{6pt}%
   /setlength{/parindent}{0pt}/setlength{/RaggedRightParindent}{0pt}/RaggedRight}
  {/end{minipage}/par/vspace{10pt}}

% The impact statement, ruled off from the abstract above it and the keywords
% below it. The same box the typst masthead draws: a half-point rule 6pt clear
% of the text on every side, with 9pt above and below.
/newsavebox{/apajouimpactbox}
/newenvironment{apajouimpact}
  {/par/vspace{3pt}%
   /setlength{/fboxrule}{0.5pt}/setlength{/fboxsep}{6pt}%
   /begin{lrbox}{/apajouimpactbox}%
   /begin{minipage}{/dimexpr/linewidth-2/fboxsep-2/fboxrule/relax}}
  {/end{minipage}/end{lrbox}%
   /noindent/fbox{/usebox{/apajouimpactbox}}/par/vspace{3pt}}

% The orcid lines of the author note. Each is one short name-and-link
% paragraph, so justifying it stretches the gaps across the measure and
% indenting it leaves the names out of line with each other. They are set
% flush left and unindented, while the prose paragraphs of the note keep the
% indent they are given, which is what the typst masthead does with them.
/newenvironment{apajouorcid}
  {/par/setlength{/parindent}{0pt}%
   /setlength{/RaggedRightParindent}{0pt}/RaggedRight}
  {/par}

% The author note, at the foot of the first page under a rule.
%
% A journal sets a short note at the foot of the first column and a long one
% across the foot of the page. apaquarto's typst format measures the note and
% does both, and so does this: the note is set once into a box the width of a
% column, and if it fits in a third of the text height it stays in that box at
% the foot of the column. A longer one is set again across the page, since a
% note that fills a column leaves the article nowhere to start.
%
% A third of the text height is the share typst measures against, and it is
% also about what latex will hold at the foot of a column: bottomfraction is
% half a page in journal mode, so a note under a third always fits and never
% drifts to a later page.
%
% The environment collects its body so that it can be set twice, which is what
% environ is for. A footnote would be the natural way to put something at the
% foot of a column, but latex sets a footnote only if it fits the column it was
% raised in and breaks it across pages otherwise, and on a page of floats that
% sent the output routine round in circles.
%
% stfloats is what allows a float spanning both columns to be set at the foot
% of a page. Latex on its own will only put one at the top, and silently drops
% a starred float asked for anywhere else.
/usepackage{stfloats}
/usepackage{environ}
/newlength{/apajounoteindent}
/setlength{/apajounoteindent}{1.1em}
/newsavebox{/apajounotebox}

/newcommand{/apajounotetext}{%
  /fontsize{9}{11}/selectfont
  /setlength{/parindent}{/apajounoteindent}%
  /setlength{/RaggedRightParindent}{/apajounoteindent}/RaggedRight}

/newcommand{/apajounoterule}{%
  /noindent/rule{/linewidth}{0.5pt}%
  /par/vspace{/dimexpr0.4em+3pt/relax}}

/NewEnviron{apajounote}{%
  /sbox{/apajounotebox}{%
    /begin{minipage}{/columnwidth}%
      /apajounotetext/BODY
    /end{minipage}}%
  /ifdim/dimexpr/ht/apajounotebox+/dp/apajounotebox/relax>/dimexpr/textheight/3/relax
    /begin{figure*}[b]%
      /apajounoterule
      /apajounotetext/BODY
    /end{figure*}%
  /else
    /begin{figure}[b]%
      /apajounoterule
      /usebox{/apajounotebox}%
    /end{figure}%
  /fi}

% --- Journal masthead -------------------------------------------------------
% The band a published article carries at the head of its first page: the
% journal's logo at the left with its name opposite, a rule under both, and the
% copyright on the left beneath it with the issue and the doi on the right.
% apaquarto's typst format sets the same band, measured off the Journal of
% Educational Psychology, and the sizes here are the ones named there: a 12pt
% name over 8pt metadata, a half-point rule, and a half-inch band the logo is
% fitted to whatever its proportions. The space around the rule and the
% leading of the metadata are not the values typst names -- see the two
% commands below, which say why.
%
% frontmatter.lua builds the masthead out of the four commands below and
% formatlatex.lua hands it to /twocolumn, which is the only thing that sets
% material across both columns of a two-column page.
/newlength{/apamastheadrulewidth}
/setlength{/apamastheadrulewidth}{0.5pt}
/newlength{/apalogoheight}
/setlength{/apalogoheight}{0.5in}
/newsavebox{/apamastheadbox}

% A journal sets its masthead above where the text of the page begins, so it is
% lifted into the top margin by half of that margin, which is what the typst
% format does. Written as half the distance from the top of the paper to the
% top of the text block, so that a document setting its own margin is followed
% rather than overruled.
%
% Half the margin above the running head would be the closer reading of what
% typst lifts by, and it puts the band within a point of typst's. It also
% makes the box that much taller, which pushes the title, the abstract and the
% first line of the body down by the same amount and puts all four of those
% out by eight points. The band is the part that can afford to be out: this
% sets it about five points higher than typst sets it, and everything under it
% level to the point.
/newcommand{/apamastheadlift}{%
  /vspace*{-/dimexpr(1in+/voffset+/topmargin+/headheight+/headsep)/2/relax}}

% Declarations rather than commands taking an argument: the metadata below is
% two lines with a break between them, which reads more plainly set into a
% paragraph than passed through a pair of braces.
/newcommand{/apamastheadname}{/fontsize{12}{14}/selectfont}
% 8pt on 10.75, which is the leading the typst masthead sets its two lines of
% metadata at. The 9.6 an 8pt size would take by the usual ratio set them a
% point and a quarter closer together than that.
/newcommand{/apamastheadsmall}{/fontsize{8}{10.75}/selectfont}

/newcommand{/apamastheadlogo}[1]{/includegraphics[height=/apalogoheight]{#1}}

% The logo and the journal's name, sitting on one baseline at either edge.
/newcommand{/apamastheadhead}[2]{%
  /noindent #1/hfill{/apamastheadname #2}/par}

% nointerlineskip on either side of the rule, and the space written out in
% full. A rule set as its own paragraph takes a whole line's leading above it
% and the metadata under it takes another, which put the rule a line lower
% than typst sets it and opened the band out by a line at each turn.
%
% The two lengths are not the 3pt the typst masthead names, because typst adds
% that to the space a grid row carries of its own and nointerlineskip here
% leaves none to add it to. They are measured instead: they put the rule and
% the first line of the metadata within a point of where the typst masthead
% puts them, which is what this band is meant to look like.
/newcommand{/apamastheadline}{%
  /par/nointerlineskip/vspace{5.5pt}%
  /noindent/rule{/linewidth}{/apamastheadrulewidth}%
  /par/nointerlineskip/vspace{6.5pt}}

% The copyright and the issn on the left, the issue and the doi on the right,
% each pair set as two lines. Half the measure each, less half the gutter the
% typst masthead leaves between them.
/newcommand{/apamastheadfoot}[2]{%
  /noindent
  /begin{minipage}[t]{/dimexpr0.5/linewidth-0.5em/relax}%
    /apamastheadsmall/raggedright #1%
  /end{minipage}/hfill
  /begin{minipage}[t]{/dimexpr0.5/linewidth-0.5em/relax}%
    /apamastheadsmall/raggedleft #2%
  /end{minipage}/par/vspace{1em}}

% The first-line indent of a published article, which is a good deal smaller
% than the half inch a manuscript takes: apa7 sets it at about 1.1 times the
% body size. Written in em so that it follows the type size, whatever the
% writer has asked for.
%
% RaggedRight reads RaggedRightParindent when it is switched on, so the length
% is set first and RaggedRight asked for again afterwards; setting parindent
% alone would be undone by the ragged setting already in force.
%
% Written into the body rather than the preamble. The half inch a manuscript
% takes is set further down the preamble than an in-header include reaches, so
% a preamble setting of this would simply be undone.
/newlength{/apajouparindent}
/setlength{/apajouparindent}{1.1em}
/newcommand{/apajournalindent}{%
  /setlength{/parindent}{/apajouparindent}%
  /setlength{/RaggedRightParindent}{/apajouparindent}%
  /RaggedRight}

% A longtable refuses to be set in two columns, and quarto writes every
% markdown table as a longtable while flextable builds its tables out of one.
% A longtable earns its keep by breaking over pages, which a table inside a
% column cannot do anyway, so in two columns each one is set as an ordinary
% tabular instead: the row and rule syntax is the same.
%
% What the conversion has to get right is the four commands that mark a
% longtable's repeating head and foot. They do not merely say where the head
% and the foot end; they say which rows repeat on every page and which are set
% once, and in the source the foot stands above the body. Made to do nothing,
% as they were, every row was set where it stood: a table with a repeating
% head printed the head twice, a pandoc table printed its /bottomrule under
% the /midrule and had no rule at the foot, and a flextable printed its note
% row between the head and the body.
%
% So the head that repeats is thrown away, the first head being the one that
% is wanted, and the foot is collected and set at the end of the table where
% it belongs. Both grabs are /long, since the rows of a flextable have blank
% lines between them, and the one that saves the foot is wrapped in /noalign:
% an assignment between two rows would otherwise open a cell, and the row
% after it lost its first column.
/makeatletter
/newtoks/apa@lt@foot
/long/def/apa@lt@firsthead#1/endhead{/apa@lt@head}
/def/apa@lt@head{/apa@lt@grab}
/long/def/apa@lt@grab#1/endlastfoot{/noalign{/global/apa@lt@foot{#1}}}
/newcommand{/apalongtableastabular}{%
  /renewenvironment{longtable}[2][]{%
    /global/apa@lt@foot{}%
    /let/endfirsthead/apa@lt@firsthead
    /let/endhead/apa@lt@head
    /let/endfoot/relax
    /let/endlastfoot/relax
    /begin{tabular}{##2}%
  }{%
    /the/apa@lt@foot
    /end{tabular}%
  }}
/makeatother

% --- Headings ---------------------------------------------------------------
% titlespacing rather than titlespacing*: the starred form takes the indent off
% the paragraph that follows a heading, which is latex's habit and not APA's.
% APA indents the first line of every paragraph, the one under a heading along
% with the rest.
% The five APA levels. One to three are their own paragraphs; four and five run
% into the text that follows them, which is why they are set as runin.
% Numbering is off: APA headings are never numbered.
/setcounter{secnumdepth}{-1}

/titleformat{/section}[block]
  {/normalfont/bfseries/centering}{}{0pt}{}
/titlespacing{/section}{0pt}{/baselineskip}{0pt}

/titleformat{/subsection}[block]
  {/normalfont/bfseries/raggedright}{}{0pt}{}
/titlespacing{/subsection}{0pt}{/baselineskip}{0pt}

/titleformat{/subsubsection}[block]
  {/normalfont/bfseries/itshape/raggedright}{}{0pt}{}
/titlespacing{/subsubsection}{0pt}{/baselineskip}{0pt}

% No [.] suffix: apaheader.lua has already put a full stop at the end of a
% level four or five heading, and only when the heading did not end in one
% already. A stop added here as well doubled every one of them, and turned a
% heading ending in a question mark into "...?.".
/titleformat{/paragraph}[runin]
  {/normalfont/bfseries}{}{0pt}{}
/titlespacing{/paragraph}{/parindent}{/baselineskip}{0.5em}

/titleformat{/subparagraph}[runin]
  {/normalfont/bfseries/itshape}{}{0pt}{}
/titlespacing{/subparagraph}{/parindent}{/baselineskip}{0.5em}

% --- Title page -------------------------------------------------------------
% The title itself, bold and centred, set a third of the way down the page in
% manuscript mode and at the top of the body otherwise.
/newcommand{/apatitle}[1]{%
  {/centering/bfseries #1/par}}

% The authors and their affiliations, centred under the title.
/newenvironment{apaauthor}
  {/par/centering}
  {/par}

% A title-page heading: Author Note, Abstract, Impact Statement. Centred and
% bold, like a level-one heading, but not part of the document's outline.
/newcommand{/apatitlepageheading}[1]{%
  {/centering/bfseries #1/par}}

% The first paragraph of the abstract and of the impact statement, which APA
% sets flush left rather than indented.
/newenvironment{apanoindent}
  {/par/setlength{/parindent}{0pt}}
  {/par}

% --- Figures and tables -----------------------------------------------------
% APA puts the number above the title, the title above the figure, and the note
% below it, all flush left. The blocks arrive already in that order, so these
% only set them.
/newcommand{/apafloattitle}[1]{%
  {/par/noindent/bfseries #1/par}}
/newcommand{/apafloatcaption}[1]{%
  {/par/noindent/itshape #1/par}}

% The blank line APA leaves between a figure's title and the figure itself.
%
% 1.2em, which is one line of text at the ordinary leading ratio and comes to
% the 14 points apa7 leaves there. Not normalbaselineskip: in a double-spaced
% document that is the stretched leading, twice as much, and it opened a gap
% ten points wider than apaquarto-pdf's.
%
% Figures only. A table follows its title immediately, and apa7 in fact leaves
% a little less room there than this format already does.
/newcommand{/apafigureskip}{/vspace{1.2em}}
/newenvironment{apafloatnote}
  {/par/setlength{/parindent}{0pt}/noindent}
  {/par}

% The air between the rule at the foot of a table and the note under it.
%
% Half an em, which is what apa7 leaves there. Without it the note's first line
% comes within a point of the rule and reads as one more row of the table.
% This was a journal-mode setting while a longtable inside a float was silently
% dropping the rule that closes it in every other mode; with that rule back,
% every mode needs the air under it.
%
% The gap is written by floatlatex.lua, which knows a table from a figure. A
% figure's note follows the picture, which needs no rule cleared.
/newlength{/apatablenotegap}
/setlength{/apatablenotegap}{0.5em}
/newcommand{/apatablenoteskip}{/par/vspace{/apatablenotegap}}

% A panel of a multipanel figure: its label above it, and its own note centred
% underneath, both set inside the panel's own column rather than across the
% figure the way the figure's note is.
%
% A minipage starts a fresh vertical list and works out its baseline skip from
% the font as it stands, which is the unstretched one: the label and the note
% came out single spaced inside a document set double. /selectfont asks for the
% skip to be worked out again, and /baselinestretch is still whatever the
% document set, so the panel is spaced like the rest of the text.
/newenvironment{apapanel}[1]
  {/begin{minipage}[t]{#1/linewidth}/selectfont/centering}
  {/end{minipage}}
/newcommand{/apapanellabel}[1]{%
  {/par/noindent #1/par/smallskip}}
/newenvironment{apapanelnote}
  {/par/smallskip/setlength{/parindent}{0pt}/centering}
  {/par}

% --- References -------------------------------------------------------------
% A hanging indent of half an inch, which is what the reference list wants and
% what quarto's csl div is given here.
/newenvironment{apareferences}
  {/par/setlength{/parindent}{0pt}%
   /everypar{/hangindent=0.5in/hangafter=1}}
  {/par}

% Quarto sets the reference list as its own CSLReferences list, which hangs its
% entries by cslhangindent and not by the environment above. That length is
% 1.5em out of the box, about two fifths of what APA asks for, so it is given
% the half inch apa7 uses. Guarded: a document with no bibliography never
% defines it.
/makeatletter
/@ifundefined{cslhangindent}{}{/setlength{/cslhangindent}{0.5in}}
/makeatother

% A published article hangs its references by the same amount it indents the
% first line of a paragraph by, which is a good deal less than the half inch a
% manuscript takes: half an inch of a column this narrow leaves the turned
% lines of an entry with little room to run in. apaquarto's typst format sets
% the hanging indent to joufirstlineindent in journal mode for the same
% reason, and the two lengths come to within a fifth of a point of each other.
% Set for jou by formatlatex.lua.
/makeatletter
/newcommand{/apajouhangindent}{%
  /@ifundefined{cslhangindent}{}{/setlength{/cslhangindent}{/apajouparindent}}}
/makeatother
/makeatletter
/@ifpackageloaded{caption}{}{/usepackage{caption}}
/AtBeginDocument{%
/ifdefined/contentsname
  /renewcommand*/contentsname{Table of contents}
/else
  /newcommand/contentsname{Table of contents}
/fi
/ifdefined/listfigurename
  /renewcommand*/listfigurename{List of Figures}
/else
  /newcommand/listfigurename{List of Figures}
/fi
/ifdefined/listtablename
  /renewcommand*/listtablename{List of Tables}
/else
  /newcommand/listtablename{List of Tables}
/fi
/ifdefined/figurename
  /renewcommand*/figurename{Figure}
/else
  /newcommand/figurename{Figure}
/fi
/ifdefined/tablename
  /renewcommand*/tablename{Table}
/else
  /newcommand/tablename{Table}
/fi
}
/@ifpackageloaded{float}{}{/usepackage{float}}
/floatstyle{ruled}
/@ifundefined{c@chapter}{/newfloat{codelisting}{h}{lop}}{/newfloat{codelisting}{h}{lop}[chapter]}
/floatname{codelisting}{Listing}
/newcommand*/listoflistings{/listof{codelisting}{List of Listings}}
/makeatother
/makeatletter
/makeatother
/makeatletter
/@ifpackageloaded{caption}{}{/usepackage{caption}}
/@ifpackageloaded{subcaption}{}{/usepackage{subcaption}}
/makeatother
/setapashorttitle{IMAGE PANELS}
/AtBeginDocument{/apalongtableastabular}
/usepackage{bookmark}
/IfFileExists{xurl.sty}{/usepackage{xurl}}{} % add URL line breaks if available
/urlstyle{same}
/hypersetup{
  colorlinks=true,
  linkcolor={blue},
  filecolor={Maroon},
  citecolor={blue},
  urlcolor={blue},
  pdfcreator={LaTeX via pandoc}}


/author{Test Author}
/date{}
/begin{document}


/hfill/break

/hfill/break

/vspace*{3/baselineskip}

/apatitle{

Multipanel Notes, Image Panels

}

/begin{apaauthor}

/hfill/break

Test Author

Test University

/end{apaauthor}

/hfill/break

/hfill/break

/apatitlepageheading{

Author Note

}

Correspondence concerning this article should be addressed to Test
Author, Test University, Email:
/href{mailto:test@example.com}{/nolinkurl{test@example.com}}

/clearpage

/apatitle{

Multipanel Notes, Image Panels

}

/section{Method}/label{method}

See Figure~/ref{fig-images}.

/begin{figure}[H]

/setcounter{figure}{0}

/refstepcounter{figure}%
/makeatletter/def/@currentlabel{1}/makeatother%
/label{fig-images}

/apafloattitle{Figure 1}

/apafloatcaption{Two Image Panels}

/apafigureskip

/begin{apapanel}{0.4900}%
/textbf{Panel A}. First panel.

/apafigureskip

/pandocbounded{/includegraphics[keepaspectratio,alt={First panel.}]{sampleimage.png}}

/begin{apapanelnote}

A note belonging to the left panel.

/end{apapanelnote}
/end{apapanel}%
/hfill
/begin{apapanel}{0.4900}%
/textbf{Panel B}. Second panel.

/apafigureskip

/pandocbounded{/includegraphics[keepaspectratio,alt={Second panel.}]{sampleimage.png}}

/begin{apapanelnote}

A note belonging to the right panel.

/end{apapanelnote}
/end{apapanel}%
/par

/begin{apafloatnote}

/emph{Note}. A~note belonging to the whole figure.

/end{apafloatnote}

/end{figure}

/section{Discussion}/label{discussion}

Text after the float.




/end{document}
